% ML Project 1 report - CS-433, EPFL, Fall 2026
% Uses the course template class (IEEEtran). If IEEEtran.cls is not found
% (e.g. outside the course template folder), it falls back to a plain
% two-column article so the file still compiles for previewing.
\IfFileExists{IEEEtran.cls}
  {\documentclass[10pt,conference,compsocconf]{IEEEtran}}
  {\documentclass[10pt,twocolumn,a4paper]{article}
   \usepackage[margin=1.8cm]{geometry}}

\usepackage{amsmath}
\usepackage{booktabs}
\usepackage{xcolor}
\usepackage[hidelinks]{hyperref}

% Placeholders for parts still to fill in
\newcommand{\todo}[1]{\textcolor{red}{[TODO: #1]}}

\begin{document}

\title{Predicting Coronary Heart Disease from Survey Data:\\
Data Cleaning Matters More than the Model}

\author{Sergi Boned Ferrer, \todo{teammate 2}, \todo{teammate 3}\\
School of Computer and Communication Sciences, EPFL, Switzerland}

\maketitle

\begin{abstract}
We predict coronary heart disease (MICHD) from 321 self-reported features of
the 2015 BRFSS survey, covering 328{,}135 respondents of whom only 8.8\% are
positive. Using only NumPy, we compare the six linear methods of the course
and an $L_2$-regularized SVM on one shared preprocessing pipeline, evaluated
with 5-fold cross-validated F1. Treating the survey's ``don't know/refused''
codes as missing values and one-hot encoding categorical answers gives the
largest improvement (+0.017 F1), whereas regularization strength, squared
features and class weighting change nothing beyond noise. Logistic regression
with an F1-tuned decision threshold is the best model (CV F1
$0.429 \pm 0.003$, test F1 0.440); the SVM stays about 0.01 below it.
\end{abstract}

\section{Introduction}

Cardiovascular diseases are among the leading causes of death worldwide, and
early identification of people at risk can support prevention. We predict
whether a respondent of the Behavioral Risk Factor Surveillance System
(BRFSS)~\cite{brfss} reports coronary heart disease from their lifestyle and
health answers, as a binary classification problem.

The data is hard for three reasons. Many questions are asked only to some
respondents, so most columns are largely empty. Answers such as ``don't know''
or ``refused'' are stored as numbers (7, 9, 77, 99, \dots), which a linear
model reads as real amounts. Finally, only 8.8\% of respondents are positive,
so accuracy is misleading and we optimize F1 instead.

Our contributions are: (i) a preprocessing pipeline that handles survey codes,
missing values and categorical answers using only training statistics;
(ii) a comparison of the six required methods and an SVM on identical data and
identical cross-validation folds; and (iii) an ablation study showing which
steps improve F1 and which do not.

\section{Models and Methods}

\subsection{Data and preprocessing}

The training set has 328{,}135 respondents and the test set 109{,}379, each
with 321 features. Labels in $\{-1,1\}$ are mapped to $\{0,1\}$. All statistics
below are computed on the training set and applied unchanged to the test set.

\begin{enumerate}
  \item \textbf{Survey codes.} If a column's training maximum is 9, the values
  7 and 9 become missing; if 99, then 77 and 99 become missing and 88
  (``none'') becomes 0; if 999, then 777 and 999 become missing and 888 becomes
  0; if 9999, then 7777 and 9999 become missing.
  \item \textbf{Column removal.} We drop 9 survey-administration columns
  (state, interview date, sequence numbers) and every column with more than
  80\% missing values, keeping 195.
  \item \textbf{One-hot encoding.} The 50 columns with 3 to 10 distinct values
  are treated as categorical. Each becomes one indicator per value plus one
  ``missing'' indicator, giving 277 binary features.
  \item \textbf{Numeric features.} The remaining 145 columns are imputed with
  the training median and standardized to zero mean and unit variance
  (z-score rather than min-max scaling, which is less sensitive to outliers).
  \item \textbf{Bias.} A constant column is added, for 423 features in total.
\end{enumerate}

\subsection{Models}

We use the six methods of the course: least squares, ridge regression
($\lambda = 10^{-4}$), linear regression with the MSE loss by gradient descent
(GD, $\gamma = 0.1$) and by stochastic gradient descent (SGD, batch size 1,
$10^5$ steps, $\gamma = 10^{-3}$), logistic regression, and $L_2$-regularized
logistic regression ($\lambda = 10^{-3}$, $\gamma = 0.5$, 500 iterations, zero
initialization). Regression models output a score; logistic models output a
probability through the sigmoid $\sigma$.

\textbf{Weighted logistic regression.} To test whether the class imbalance
should be handled in the loss, positive examples get weight
$w_+ \in \{3, 10\}$ and negatives weight 1, with gradient
\begin{equation}
\nabla L(w) = \frac{1}{\sum_n c_n} X^\top \big(c \odot (\sigma(Xw) - y)\big)
+ 2\lambda w,
\end{equation}
where $c_n = w_+$ if $y_n = 1$ and $c_n = 1$ otherwise.

\textbf{$L_2$-regularized SVM.} With labels $y_n \in \{-1, 1\}$, we minimize
$\frac{1}{\sum_n c_n}\sum_n c_n \max(0, 1 - y_n x_n^\top w) +
\frac{\lambda}{2}\|w\|^2$ with $\lambda = 10^{-3}$, using full-batch
subgradient descent (500 steps, step size $0.1/\sqrt{t+1}$) and returning the
average of all iterates. Without averaging, the last iterate oscillated and
the F1 varied from 0.23 to 0.42 across folds.

\subsection{Evaluation protocol}

Models are compared with 5-fold cross-validation (random seed 1, the same
folds for every model), reporting the mean $\pm$ standard deviation of F1. In
each fold the decision threshold is chosen on the validation fold by a grid
search with step 0.01 (0.05--0.60 for probabilities, 0.05--0.95 for weighted
models, $-2$ to 2 for SVM scores). The ablation steps use a single 80/20 split
with the same seed. The final model is retrained on all training data and
applied to the test set with the mean of the five fold thresholds.

\section{Results}

Table~\ref{tab:ablation} shows the ablation: each row adds one change to the
row above. The two preprocessing steps give the largest gains (+0.013 and
+0.004), while tuning $\lambda$ and adding squared features give none.

\begin{table}[htbp]
\centering
\caption{Ablation on a single 80/20 split (validation F1), with test F1
where submitted.}
\label{tab:ablation}
\begin{tabular}{@{}clcc@{}}
\toprule
Step & Change & Val.\ F1 & Test F1 \\
\midrule
1 & Ridge, basic cleaning & 0.4106 & \\
2 & Reg.\ logistic, not ridge & 0.4136 & \\
3 & + survey codes as missing & 0.4269 & \\
4 & + one-hot encoding & 0.4308 & \\
5 & $\lambda$ tuned, best $10^{-3}$ & 0.4315 & 0.439 \\
6 & + squared numeric features & 0.4315 & \\
\bottomrule
\end{tabular}
\end{table}

Table~\ref{tab:methods} compares all methods on the final preprocessing with
5-fold CV. Logistic regression is best, with or without regularization and
class weights; ridge and the weighted SVM are about 0.01 lower, three standard
deviations. The selected threshold is stable across folds (0.19--0.21 for
unweighted logistic), and both submitted models score about 0.01 higher on
the test set than in CV, so CV ranks the models reliably.

\begin{table}[htbp]
\centering
\caption{Method comparison, 5-fold CV F1 (mean $\pm$ std).}
\label{tab:methods}
\begin{tabular}{@{}lcc@{}}
\toprule
Method & CV F1 & Test F1 \\
\midrule
Weighted logistic, $w_+ = 3$ & $0.4287 \pm 0.0033$ & 0.440 \\
Logistic regression & $0.4285 \pm 0.0029$ & \\
Reg.\ logistic, $\lambda = 10^{-3}$ & $0.4281 \pm 0.0035$ & \\
Weighted logistic, $w_+ = 10$ & $0.4276 \pm 0.0030$ & \\
Ridge, $\lambda = 10^{-4}$ & $0.4197 \pm 0.0033$ & \\
SVM, $w_+ = 10$ & $0.4193 \pm 0.0027$ & \\
SVM, $w_+ = 5$ & $0.4166 \pm 0.0047$ & \\
SVM, unweighted & $0.4027 \pm 0.0020$ & \\
MSE SGD & $0.147 \pm 0.020$ & \\
MSE GD & diverged & \\
Least squares & singular & \\
\bottomrule
\end{tabular}
\end{table}

\section{Discussion}

\textbf{Preprocessing drives performance.} Survey codes such as 99
(``refused'') are far larger than real answers, so a linear model gives them a
large, meaningless effect; turning them into missing values gave the largest
single gain. One-hot encoding helps for the same reason: categorical answers
stop being treated as amounts, and ``don't know'' becomes its own indicator.

\textbf{Regularization does not matter here.} With 328{,}135 examples and 423
features the model is far from overfitting, so $\lambda$ between $10^{-5}$ and
$10^{-2}$ gives the same F1, for logistic regression and SVM alike. Squared
features did not help either, likely because many numeric columns are already
grouped (e.g.\ 5-year age bands). 500 and 1000 iterations give the same
training loss, so the logistic models have converged.

\textbf{Threshold tuning already handles the imbalance.} For logistic
regression, class weighting shifts the predicted probabilities upward and the
selected threshold moves with them (0.21, 0.44 and 0.52 for $w_+ = 1, 3, 10$),
leaving F1 unchanged. The test set confirms this: the unweighted and the
$w_+ = 3$ models score 0.439 and 0.440, a difference within noise.

\textbf{Hinge loss ranks worse.} The unweighted SVM pushes most negatives to a
score of exactly $-1$, where the hinge loss stops penalizing them; the best
threshold sits just below $-1$ in every fold and many respondents share nearly
the same score, which limits threshold tuning. Weighting positives spreads the
scores and recovers most of the gap (0.403 to 0.419), but the logistic loss,
which grades every example, remains better for an F1-tuned threshold.

\textbf{Method notes.} Least squares cannot be solved on this feature matrix:
the indicators of each categorical variable sum to the bias column, so
$X^\top X$ is singular; the $\lambda I$ term of ridge makes it invertible. MSE
gradient descent diverges with $\gamma = 0.1$, and SGD with batch size 1 is
far from converged after $10^5$ steps (less than half a pass over the data).

\textbf{Limitations.} The survey-code rule is a heuristic based on each
column's maximum and misses some cases (e.g.\ 8 = ``never'' in a few
questions). The ablation uses a single split, whose noise (about 0.003) is
close to some of the differences. All our models are linear.

\section{Summary}

On the BRFSS heart-disease task, careful preprocessing (survey codes as
missing, one-hot encoding, median imputation, standardization) raised the
validation F1 from 0.411 to 0.431, while regularization, squared features,
class weights and the hinge loss gave no further gain. Our submitted model is
weighted logistic regression ($w_+ = 3$, $\lambda = 10^{-3}$) with an F1-tuned
threshold, at CV F1 $0.429 \pm 0.003$ and test F1 0.440; \texttt{run.py}
reproduces this submission exactly.

\begin{thebibliography}{1}
\bibitem{brfss}
Centers for Disease Control and Prevention, ``Behavioral Risk Factor
Surveillance System, 2015 annual data,''
\url{https://www.cdc.gov/brfss/annual_data/annual_2015.html}.
\end{thebibliography}

\end{document}
